% !TEX program = pdflatex
% !TEX encoding = UTF-8 Unicode

% ====================================================================
% ARCHIVO PRINCIPAL — Práctica Grupal Temas 4 y 5
% Asignatura: Dirección de Recursos Humanos I
% Doble Grado ADE-Ingenierías — Universidad de Granada
% Empresa analizada: CaixaBank
%
% INSTRUCCIONES DE COMPILACIÓN:
%   pdflatex practicat4t5.tex && pdflatex practicat4t5.tex
%   (dos pasadas para generar índice y referencias correctamente)
%
% Para la bibliografía, si añadís citas al archivo library.bib:
%   pdflatex practicat4t5.tex && bibtex practicat4t5 && pdflatex practicat4t5.tex && pdflatex practicat4t5.tex
% ====================================================================

\documentclass[print, color]{ugrTFG}

% -------------------------------------------------------------------
%  VERSIÓN ELECTRÓNICA PARA TABLETA
%  Cambiad "print" por "tablet" para generar un PDF adaptado a tabletas.
% -------------------------------------------------------------------

% -------------------------------------------------------------------
%  INFORMACIÓN DEL TRABAJO Y LOS AUTORES
% -------------------------------------------------------------------

\renewcommand{\miTipoTrabajo}{Práctica Grupal}
\renewcommand{\miAsignatura}{Dirección de Recursos Humanos I}

\newcommand{\miTitulo}{Formación y Desarrollo, Evaluación y Gestión del Rendimiento y Retribución de la Empresa \xspace}

% Para un trabajo grupal, \miNombre contiene todos los autores.
% En la portada aparecerá como "Presentado por:" seguido de los nombres.
% Ajustadlos con vuestros nombres reales.
\newcommand{\miNombre}{%
  José Angel Carretero Montes \\ 
  David Bacas Posadas \\
  Julian Carrion Tovar \\
  Jesús Rodriguez Gonzalez \\
  Ismael Sallami Moreno \\
}

\newcommand{\miGrado}{Doble Grado ADE-Ingenierías}
\newcommand{\miFacultad}{Facultad de Ciencias Empresariales y Escuela Técnica Superior de Ingenierías}
\newcommand{\miUniversidad}{Universidad de Granada}

% Nombre del profesor/a responsable de la asignatura
\newcommand{\miTutor}{
  Javier Aguilera Caracuel \\ \emph{Departamento de Organización de Empresas II}
}

\newcommand{\miCurso}{2025-2026}

% -------------------------------------------------------------------
%  PAQUETES ADICIONALES PARA ESTE TRABAJO
% -------------------------------------------------------------------

\usepackage{tabularx}       % Tablas con columnas de anchura automática
\usepackage{multirow}       % Celdas que ocupan varias filas
\usepackage{float}          % Control de posición de figuras [H]
\usepackage{eurosym}        % Símbolo del euro: \euro{}
\usepackage{enumitem}       % Control avanzado de listas
\usepackage{tikz}           % Gráficos vectoriales (organigrama)
\usepackage{rotating}       % Figuras apaisadas (sidewaysfigure, sin problemas twoside)

% Cajas de color para recordatorios al equipo
\usepackage[most]{tcolorbox}
\tcbuselibrary{skins, breakable}

% Color corporativo UGR
\definecolor{ugrazul}{RGB}{0,75,135}

% Entorno "notaequipo": caja amarilla para recordatorios y preguntas guía.
% Usadlo cuando queráis dejar instrucciones a vuestros compañeros sobre
% qué información hay que buscar o qué preguntar en la entrevista.
\newtcolorbox{notaequipo}{
  enhanced,
  breakable,
  colback=yellow!15!white,
  colframe=orange!70!black,
  fonttitle=\sffamily\bfseries,
  title={Recordatorio del equipo},
  left=6pt, right=6pt, top=4pt, bottom=4pt,
  before skip=8pt, after skip=8pt
}

% -------------------------------------------------------------------
%  METADATOS PDF
% -------------------------------------------------------------------
\hypersetup{
  pdftitle={\miTitulo},
  pdfauthor={Grupo --- Dirección de Recursos Humanos I --- UGR},
  pdfsubject={Práctica Grupal — CaixaBank}
}

% ====================================================================
\begin{document}
% ====================================================================

\maketitle

% -------------------------------------------------------------------
%  FRONTMATTER
%  (páginas previas numeradas en romano: resumen, índice)
% -------------------------------------------------------------------
\frontmatter

% Resumen ejecutivo del trabajo
% !TeX root = ../practicat6t7t8.tex
% !TeX encoding = utf8

\chapter{Resumen}

El presente trabajo analiza las políticas de Formación y Desarrollo
(Tema~6), Evaluación y Gestión del Rendimiento (Tema~7) y Retribución
(Tema~8) aplicadas en CaixaBank, S.A., primer grupo bancario español
por volumen de activos, con más de 45.000 empleados.

En el ámbito de la formación, se examina el ciclo completo del programa
formativo de la entidad —identificación de necesidades, diseño,
implantación y evaluación—, así como el tiempo y coste asociados.
Se analizan igualmente las nuevas tendencias adoptadas: la plataforma
de \textit{e-learning} Virtaula, el modelo 70-20-10, la estructura de
universidades corporativas (escuelas especializadas internas),
el \textit{outdoor training} y la gamificación. Respecto a la gestión
de la carrera profesional, se describe el proceso de planificación
en sus cuatro fases —premisas, valoración, dirección y desarrollo—,
detallando las herramientas de rotación de puestos, \textit{coaching},
\textit{mentoring} e \textit{mentoring} inverso empleadas por la entidad,
junto con las tendencias de \textit{networking}, empleabilidad y marca
personal.

En el ámbito de la evaluación del rendimiento, se estudia la delimitación
de la evaluación (qué, cuándo y quién evalúa), los métodos objetivos,
subjetivos y mixtos utilizados, así como las estrategias de reforzamiento
positivo, negativo y la entrevista de evaluación. Se presta especial
atención a las tendencias de evaluación continua y por compromisos que
CaixaBank está implantando en sustitución de los sistemas anuales
tradicionales.

Finalmente, en materia retributiva, se analiza la estructura salarial de
cinco puestos representativos con referencia al Convenio Colectivo
sectorial de Banca, los tres tipos de equidad (interna, externa e
individual), el modelo híbrido de retribución por puesto y por
competencias, los principales incentivos financieros y elementos no
financieros, y las tendencias de retribución flexible y salario emocional.

La metodología combina el análisis de fuentes documentales públicas
—memoria anual, informe de sostenibilidad y Plan Estratégico
2025--2027 de CaixaBank— con una entrevista semiestructurada a una
directora de sucursal de la entidad.

\medskip

\textbf{Palabras clave:} formación corporativa, \textit{e-learning},
gestión de carrera profesional, \textit{coaching}, \textit{mentoring},
evaluación del rendimiento, retribución flexible, salario emocional,
CaixaBank, convenio colectivo de banca.

\cleardoublepage
\endinput


% Tabla de contenidos
% !TeX root = ../tfg.tex
% !TeX encoding = utf8
%
% ====================================================================
%  TABLA DE CONTENIDOS
% ====================================================================

\phantomsection
\pdfbookmark[0]{\contentsname}{toc}

\setcounter{tocdepth}{2}      % El índice muestra hasta subsecciones
\setcounter{secnumdepth}{3}   % Se numeran hasta subsubsecciones

\tableofcontents

% Listas de figuras y tablas (descomentad si las necesitáis)
% \phantomsection
% \listoffigures

% \phantomsection
% \listoftables

\cleardoublepage


% -------------------------------------------------------------------
%  MAINMATTER
%  (cuerpo del trabajo, numerado en arábigo)
% -------------------------------------------------------------------
\mainmatter

\chapter{Formación y Desarrollo}

\section{Programa de formación utilizado (identificación, diseño e implantación)}

La gestión del aprendizaje en la organización se desarrolla siguiendo un ciclo técnico que busca transformar las capacidades de los empleados para alinearlas con la competitividad del sector bancario. Este proceso se divide en las fases que se detallan a continuación.  

\section{Programa de formación utilizado (identificación, diseño e implantación)}

En CaixaBank, la mejora continua de su plantilla se estructura partiendo de una distinción teórica fundamental que la entidad aplica en su día a día: la diferencia entre formación y desarrollo. Por un lado, la entidad utiliza la formación como un proceso reactivo centrado en el trabajo actual; por ejemplo, instruyendo a los empleados provenientes de la fusión con Bankia en el uso del software propio de la entidad o actualizando conocimientos normativos inmediatos. Por otro lado, aplica el desarrollo como un proceso proactivo y a largo plazo, preparando a empleados talentosos mediante planes de carrera para que asuman futuras posiciones de liderazgo o perfiles altamente tecnificados.

Para llevar a cabo ambos procesos, el Área de Personas y Organización de CaixaBank sigue un ciclo técnico estructurado en las siguientes fases:

\subsection{Identificación de necesidades de formación}

Esta etapa inicial busca diagnosticar las carencias del personal para justificar la inversión en aprendizaje. Siguiendo los principios teóricos, los responsables de recursos humanos de CaixaBank deben dar respuesta a tres preguntas fundamentales antes de planificar cualquier acción:

\begin{itemize}
    \item ¿Qué conocimientos y habilidades concretas en el desempeño del trabajo mejorarían la productividad laboral?
    \item ¿Cuál es el público objetivo?
    \item ¿Puedo desarrollar objetivos mensurables para esos conocimientos y habilidades?
\end{itemize}

Para responder a estas cuestiones con precisión, la entidad evalúa sus necesidades estratificando la información en tres niveles:

\begin{itemize}
    \item Nivel de organización: Se evalúan las necesidades derivadas de la estrategia global del banco. Por ejemplo, observando el Plan Estratégico 2025-2027, la dirección ha identificado que toda la organización necesita formarse urgentemente en inteligencia artificial generativa, finanzas sostenibles y atención a los nuevos ecosistemas vinculados al envejecimiento poblacional (segmento sénior).
    \item Nivel de tareas: Apoyándose en el departamento de Secretaría Técnica de Negocios y consultoras como Deloitte, se analiza qué exige cada puesto concreto. Así, detectan que el puesto de Gestor de Banca Privada requiere obligatoriamente habilidades analíticas complejas y certificaciones financieras reguladas.
    \item Nivel de individuo: Utilizando su Sistema de Información de Recursos Humanos (la plataforma en la nube SAP SuccessFactors), cruzan los datos de la evaluación del desempeño de cada trabajador con los requisitos de su puesto. Esto permite identificar de forma automatizada las brechas entre el rendimiento actual del empleado y su potencial, generando alarmas sobre qué formación específica necesita cada individuo.
\end{itemize}

\subsection{Diseño de los programas de formación}

Una vez que SAP SuccessFactors y los directores de área han detectado las necesidades, se procede al diseño del programa. En esta etapa, la teoría exige concretar los siguientes parámetros:

\begin{itemize}
    \item ¿Qué tipo de formación impartir?
    \item ¿Qué grado de aprendizaje alcanzar?
    \item ¿Para quién?
    \item ¿Quién la va a impartir?
\end{itemize}

CaixaBank diseña los contenidos estructurándolos en diferentes grados de aprendizaje. Para niveles iniciales, diseñan programas de acogida (onboarding) centrados en el saber; para niveles medios, diseñan formaciones en habilidades comerciales y trabajo en equipo orientadas al poder hacer; y para el desarrollo directivo, se enfocan en habilidades conceptuales.

Respecto a quién imparte la formación, la entidad utiliza un enfoque mixto. Como formadores internos, CaixaBank emplea a sus propios supervisores y a empleados senior mediante programas de mentoría, algo esencial para transmitir la cultura organizativa y orientar a los nuevos. Como formadores externos, cuando el grado de aprendizaje requerido es muy elevado y técnico, la entidad externaliza el servicio. Por ejemplo, se apoya en alianzas con la UPF Barcelona School of Management y otras instituciones académicas para impartir los programas que otorgan a sus empleados las certificaciones europeas MiFID II o EFPA, obligatorias para el asesoramiento financiero.

\subsection{Implantación de los programas de formación}

La última fase operativa requiere decidir el entorno y los canales tecnológicos que se utilizarán, dando respuesta a dos cuestiones clave:

\begin{itemize}
    \item ¿Dónde se realizará la formación?
    \item ¿Qué medios se utilizarán?
\end{itemize}

Dada la inmensa capilaridad de su red de oficinas (miles de sucursales en todo el territorio), CaixaBank combina metodologías para asegurar que la formación llegue a toda la plantilla de manera eficiente:

\begin{itemize}
    \item Métodos en el puesto de trabajo: Se aplican mediante la rotación de puestos (moving interno) y el aprendizaje directo. Por ejemplo, la asignación de un compañero experimentado (buddy) a los nuevos gestores comerciales permite que el trabajador adquiera la experiencia resolviendo incidencias reales con clientes, lo que asegura una transferencia directa del conocimiento sin abandonar la sucursal.
    \item Métodos fuera del lugar de trabajo (y entornos virtuales): La principal herramienta de implantación masiva de CaixaBank es su plataforma corporativa de aprendizaje denominada Virtaula. Se trata de un Learning Management System (LMS) donde la entidad aloja todo su e-learning. Dentro de Virtaula, el conocimiento está segmentado en escuelas especializadas (como la Escuela de Riesgos o la Escuela de Liderazgo). 
\end{itemize}

Además, en la implantación actual, CaixaBank está utilizando licencias de Microsoft Copilot integradas en sus sistemas. Esta inteligencia artificial actúa en la fase de implantación recomendando a cada empleado, de forma personalizada, qué módulos o píldoras formativas de Virtaula debe cursar a continuación según su progreso, permitiendo un aprendizaje a ritmo individual que reduce drásticamente los costes logísticos frente a los antiguos cursos presenciales.  

\section{Evaluación de los programas formativos}

La fase de evaluación constituye el cierre del ciclo de formación en CaixaBank y tiene como propósito fundamental determinar si las actividades realizadas han satisfecho los objetivos establecidos inicialmente. Este proceso no se limita a una comprobación administrativa del cumplimiento de los cursos, sino que busca medir si la inversión ha servido para mejorar el rendimiento presente y futuro de los empleados de la entidad.

Para que este análisis sea riguroso, el departamento de recursos humanos de la organización aplica un modelo que busca dar respuesta a los siguientes niveles de control definidos en la teoría:

\begin{itemize}
\item Reacción: ¿indican las reacciones iniciales de las personas ante la experiencia formativa que el aprendizaje es relevante e inmediatamente aplicable?
\item Aprendizaje: ¿han adquirido los conocimientos y habilidades planificadas?
\item Conducta: ¿se aplica lo aprendido en el trabajo?
\item Rendimiento: ¿qué resultados se aprecian en el negocio?
\item Impacto: ¿cuál es el coste-beneficio para la empresa?
\end{itemize}

Para llevar a la práctica estos niveles de evaluación en una red tan extensa, CaixaBank utiliza diversas herramientas tecnológicas y metodológicas integradas en su ecosistema digital:

Para medir la reacción y el aprendizaje, la entidad aprovecha las capacidades analíticas de su plataforma Virtaula. Al finalizar cada módulo formativo, especialmente en los itinerarios de e-learning, el sistema requiere que el empleado complete cuestionarios de satisfacción donde se evalúa la utilidad percibida de los contenidos. Asimismo, el aprendizaje se valida mediante test y pruebas de conocimientos integradas en el propio entorno virtual. Estas pruebas son obligatorias en ámbitos de cumplimiento normativo, como la formación en blanqueo de capitales, o para obtener las certificaciones financieras externas necesarias para operar en banca privada.

La evaluación de la conducta se realiza de manera más cualitativa y a medio plazo a través de la gestión del rendimiento. La organización lleva a cabo seguimientos periódicos, usualmente a los seis meses y al año de una formación relevante o de un cambio de puesto. En estas etapas, se analiza si el trabajador aplica efectivamente lo aprendido para agilizar procesos o mejorar la atención al cliente. En este nivel, la entrevista de evaluación con el responsable directo es la herramienta clave para confirmar si las nuevas habilidades técnicas o actitudes se han incorporado con éxito al día a día de la sucursal.

Finalmente, el rendimiento y el impacto se miden contrastando los resultados de la formación con las metas estratégicas del banco. Utilizando People Analytics dentro de su sistema SAP, CaixaBank monitoriza si los programas de reskilling y reciclaje profesional han contribuido a alcanzar indicadores de negocio como el aumento del margen de intereses o la mejora de la eficiencia operativa proyectada en el Plan Estratégico 2025-2027. Este análisis final permite a la dirección determinar si el coste de los programas ha generado un beneficio tangible que justifique el presupuesto invertido en el capital humano.

\section{Análisis de tiempo y coste de los programas}

En CaixaBank, la gestión del tiempo y el coste de la formación no se analiza de forma aislada, sino como una pieza fundamental del proceso de dimensionamiento que determina la estructura global de costes del banco. La entidad entiende que cualquier actividad formativa conlleva una inversión de recursos que debe ser optimizada para asegurar que el aprendizaje sea eficiente y no interfiera excesivamente en la operatividad diaria de las oficinas.

\subsection{Gestión y optimización del tiempo formativo}

La variable tiempo es crítica en el sector bancario, donde la carga de trabajo operativa es elevada. Según la teoría, los métodos de formación en el puesto de trabajo, aunque efectivos, suelen requerir mucho tiempo y pueden interferir con el rendimiento real. CaixaBank gestiona este riesgo mediante las siguientes estrategias:

\begin{itemize}
    \item Uso de Virtaula para el aprendizaje asíncrono: Al emplear un sistema de formación por ordenador y vídeo interactivo, la entidad permite el aprendizaje a ritmo individual. Esto reduce el tiempo que los empleados deben pasar fuera de sus puestos en cursos presenciales y elimina los tiempos de desplazamiento.
    \item Microaprendizaje o Instant Learning: La plataforma permite consumir soluciones instantáneas a carencias profesionales mediante píldoras informativas breves. Esto responde a la necesidad de formación inmediata sin comprometer largas jornadas laborales.
    \item Planificación estratégica del tiempo: La entidad utiliza métricas de seguimiento para monitorizar el tiempo de contratación y el tiempo de integración de nuevos empleados. El objetivo es que la curva de aprendizaje sea lo más corta posible para que el profesional alcance su nivel de productividad óptimo en el menor plazo de tiempo.
\end{itemize}

\subsection{Análisis y estructura de costes de la formación}

Aunque la organización no publica una cifra cerrada dedicada exclusivamente a la formación, su estrategia de costes se integra en la inversión masiva de más de 5.000 millones de euros destinada a tecnología y transformación para el periodo 2025-2027. La entidad diferencia entre diversos tipos de costes según la metodología empleada:

\begin{itemize}
    \item Costes de formación interna: Son los derivados del uso de supervisores y compañeros como formadores. Aunque estos métodos son teóricamente más baratos por no requerir instalaciones externas, CaixaBank asume el coste indirecto del tiempo que estos profesionales dejan de dedicar a sus tareas comerciales para actuar como mentores.
    \item Costes de formación externa y especializada: La entidad presupuesta partidas específicas para colaborar con asesores externos y universidades. Estos programas, como los necesarios para las certificaciones financieras, tienen un coste directo elevado pero garantizan un nivel de aprendizaje superior y el cumplimiento de la normativa legal.
    \item Inversión en infraestructuras digitales: El desarrollo y mantenimiento de sistemas como SAP SuccessFactors y Virtaula supone un coste de preparación muy alto. Sin embargo, una vez implantados, permiten formar a grandes grupos de personas (los más de 45.000 empleados del grupo) de forma mucho más barata y accesible que los cursos reglados tradicionales.
\end{itemize}

Finalmente, el banco utiliza herramientas de Learning Analytics para evaluar el impacto económico. Esto permite identificar patrones de aprendizaje y adaptar los itinerarios, asegurando que cada euro invertido en formación genere un retorno medible en términos de mejora de capacidades y cumplimiento de los objetivos estratégicos del grupo.

\section{Tendencias en formación: E-learning, Método 70-20-10 y Universidades Corporativas}

CaixaBank ha transformado su modelo de aprendizaje adoptando tendencias que permiten una formación más flexible, colaborativa y alineada con los objetivos de desarrollo a largo plazo. La entidad no solo busca corregir carencias inmediatas a través de cursos tradicionales, sino que utiliza metodologías avanzadas para potenciar el talento de forma continua.

\subsection{E-learning y entornos virtuales de aprendizaje}

La formación en línea es el eje central de la estrategia de aprendizaje de la entidad. Basada en el concepto de aprendizaje 3.0, la organización fomenta un entorno autónomo y colaborativo que elimina las barreras geográficas y temporales. La herramienta principal para esta tendencia es Virtaula, el Learning Management System del banco. Esta plataforma permite la gestión de contenidos digitales y facilita el microaprendizaje o instant learning, proporcionando soluciones rápidas a necesidades profesionales puntuales.

Además, la entidad está incorporando analítica de aprendizaje (learning analytics) y big data para identificar patrones de comportamiento en los estilos de estudio de sus empleados. Esto permite crear itinerarios de aprendizaje hiper-personalizados en la nube. La reciente integración de inteligencia artificial generativa, como Microsoft Copilot, refuerza esta tendencia al asistir a los empleados en la búsqueda de contenidos relevantes dentro de la plataforma y sugerir formaciones de manera proactiva según el perfil técnico o comercial del usuario.

\subsection{El modelo de aprendizaje 70-20-10}

CaixaBank aplica de manera práctica la metodología 70-20-10, la cual establece que el aprendizaje efectivo se produce a través de tres vías diferenciadas:

En primer lugar, el 70 por ciento del aprendizaje es experiencial. Se obtiene mediante la práctica diaria, la resolución de problemas en la oficina y la participación en nuevos proyectos o tareas. La entidad potencia este bloque a través de la rotación de puestos, permitiendo que el trabajador asuma retos en diferentes áreas para adquirir una visión global del negocio.

En segundo lugar, el 20 por ciento corresponde a la interacción social. Para cubrir este porcentaje, la organización utiliza sistemas de coaching y mentoring. Un ejemplo claro es la asignación de un compañero o mentor (buddy) a las nuevas incorporaciones, facilitando que el conocimiento se transmita de forma orgánica a través del debate y el trabajo en equipo dentro de la propia sucursal.

En tercer lugar, el 10 por ciento restante es el aprendizaje formal. Este se canaliza principalmente a través de Virtaula mediante cursos, talleres y seminarios técnicos. Aunque es el porcentaje menor, es fundamental para asentar las bases teóricas y cumplir con las certificaciones normativas obligatorias del sector bancario.

\subsection{Estructura de Universidad Corporativa}

Aunque la entidad no utiliza siempre el término comercial de universidad corporativa, su estructura interna de escuelas especializadas funciona bajo este concepto teórico. Se trata de estructuras diseñadas para asegurar que el aprendizaje esté directamente conectado con la estrategia y los procesos del banco.

CaixaBank cuenta con diversas escuelas de conocimiento, como la Escuela de Riesgos, la Escuela de Liderazgo o la Escuela Comercial. Estas unidades son las encargadas de transmitir no solo los conocimientos técnicos necesarios para el sector, sino también los valores, la cultura y la historia de la organización. Este modelo asegura que la formación sea coherente en todo el grupo y que los programas de desarrollo preparen a los mandos medios y superiores para los desafíos estratégicos del futuro.

\subsection{Gamificación y nuevas dinámicas}

La organización está implementando estrategias de juego en entornos profesionales para aumentar el compromiso y reforzar habilidades como la creatividad o la comunicación interna. A través de dinámicas interactivas y test en formato de retos, la entidad busca que la formación continua tenga un impacto positivo en la motivación de los trabajadores. Estas técnicas se utilizan tanto en los procesos de integración de nuevos empleados como en los programas de actualización de competencias para perfiles directivos, facilitando un aprendizaje más dinámico y efectivo.
\section{Tendencias: Outdoor training y Gamificación.} 
\section{Planificación de carrera (Premisas).}
\section{Fase de valoración y Fase de dirección de carrera.}
\section{Fase de desarrollo (rotación, coaching, mentoring, mentoring inverso).}
\section{Tendencias en gestión de carrera (Networking, Plan de empleabilidad y Marca personal).}

\chapter{Evaluación y Gestión del Rendimiento}

\section{Delimitación de la evaluación del rendimiento (qué, cuándo y quién).}
\section{Métodos de evaluación (objetivos, subjetivos y mixtos).}
\section{Gestión y mejora del rendimiento.  }
\section{Estrategias de reforzamiento (positivo, negativo y radicales).}
\section{Entrevista de evaluación.} 
\section{Tendencias: Evaluación continua y Evaluación por compromisos.}  

\chapter{Retribución en la Empresa}
\section{Tipos de retribución (financiera/no financiera, directa/indirecta, fija/variable, prestaciones).}
\chapter{Retribución en CaixaBank}

\section{Análisis de puestos representativos y consulta del Convenio Colectivo sectorial}

La política retributiva de CaixaBank se enmarca en el \textbf{IV Convenio Colectivo del Sector de Banca}, negociado entre la Asociación Española de Banca (AEB) y los principales sindicatos del sector (CCOO, UGT y FINE). Este convenio establece las categorías profesionales, las tablas salariales mínimas y las condiciones laborales de aplicación general para todas las entidades bancarias adheridas, actuando como suelo retributivo que CaixaBank respeta y, en la mayoría de los casos, supera.

Con el fin de reflejar la diversidad de funciones y niveles jerárquicos de la entidad, se han seleccionado cinco puestos representativos distribuidos a lo largo de su estructura organizativa, desde los perfiles operativos hasta la alta dirección:

\begin{itemize}
    \item \textbf{Operativo de caja / Auxiliar administrativo (Nivel I):} Es el puesto de entrada en la estructura de sucursal. Sus funciones abarcan la atención básica al cliente en ventanilla, la gestión de efectivo y las operaciones de rutina. Según el convenio colectivo sectorial vigente, el salario base mínimo para esta categoría se sitúa en torno a los \textbf{17.500\,€ brutos anuales}, aunque en CaixaBank la retribución total efectiva, incluyendo complementos de antigüedad y pluses de convenio, alcanza habitualmente los \textbf{20.000--22.000\,€ brutos anuales}.

    \item \textbf{Gestor de Banca Personal (Nivel II/III):} Perfil comercial encargado de asesorar a clientes particulares en productos de ahorro, inversión básica y financiación. Requiere la certificación MiFID II. El convenio fija su retribución mínima en aproximadamente \textbf{22.000\,€ brutos anuales}; en CaixaBank, el paquete retributivo total —incluyendo variable comercial— oscila entre los \textbf{28.000\,€ y los 35.000\,€ brutos anuales} en función de los objetivos alcanzados.

    \item \textbf{Gestor de Banca Privada / Analista Senior (Nivel III/IV):} Especialista en asesoramiento patrimonial a clientes de alto patrimonio neto. Exige la certificación EFPA o EFA y habilidades analíticas avanzadas. Es uno de los perfiles más valorados del sector, con una retribución que en CaixaBank se aproxima a los \textbf{45.000--60.000\,€ brutos anuales}, a lo que se añade una parte variable significativa vinculada al volumen de patrimonio gestionado.

    \item \textbf{Director/a de Sucursal (Nivel IV/V):} Responsable de la gestión integral de una oficina bancaria: liderazgo del equipo, captación de negocio, cumplimiento normativo y relación con clientes clave. Su retribución en CaixaBank ronda los \textbf{50.000--70.000\,€ brutos anuales}, incluyendo un variable ligado a los resultados globales de la sucursal.

    \item \textbf{Director/a de Área o Ejecutivo/a Senior (Nivel V/Directivo):} Perfiles de alta responsabilidad que supervisan varias sucursales o lideran departamentos corporativos. Su retribución total puede superar los \textbf{100.000\,€ brutos anuales}, complementada con incentivos a largo plazo como acciones del banco o planes de co-inversión.
\end{itemize}

Cabe destacar que CaixaBank cotiza al Ibex 35 y publica anualmente su brecha salarial de género y sus ratios retributivos en el informe de sostenibilidad, donde se recoge que la retribución media de su plantilla supera en un 30--40\,\% los mínimos establecidos por el convenio sectorial, lo que refleja una apuesta decidida por la competitividad retributiva como palanca de atracción y retención del talento.

\section{Gestión de la retribución: los tres tipos de equidad}

Una gestión retributiva eficaz debe garantizar que los empleados perciban su salario como justo, tanto en relación con sus compañeros como con el mercado externo y con su propia aportación individual. CaixaBank articula su política de compensación en torno a los tres tipos clásicos de equidad retributiva.

\subsection{Equidad interna}

La equidad interna hace referencia a la coherencia salarial entre los distintos puestos de trabajo dentro de la propia organización. Para garantizarla, CaixaBank estructura su plantilla en un sistema de \textbf{niveles o bandas salariales} —denominados internamente grupos profesionales— que clasifica todos los puestos según su valor relativo para la organización, atendiendo a criterios como la responsabilidad asumida, la complejidad de las tareas, el impacto en el negocio y las competencias técnicas requeridas.

Este modelo asegura que dos empleados que desempeñen funciones de igual valor y complejidad no experimenten diferencias salariales injustificadas. La implantación de \textbf{SAP SuccessFactors} como plataforma central de gestión de personas facilita la administración de estas bandas y permite a los responsables de Recursos Humanos detectar desviaciones internas con rapidez, lo que contribuye a mantener la coherencia de la estructura retributiva en los más de 45.000 empleados del grupo.

\subsection{Equidad externa}

La equidad externa mide la competitividad de la retribución de la empresa respecto a otras organizaciones del sector o del mercado de referencia. Para garantizarla, el Área de Personas y Organización de CaixaBank realiza anualmente \textbf{estudios de benchmarking salarial} apoyándose en las encuestas de compensación de consultoras especializadas como Mercer o Willis Towers Watson, que recogen los datos del sector financiero español e internacional.

Estos estudios permiten a la entidad conocer en qué posición se sitúa su paquete retributivo en comparación con sus competidores directos —BBVA, Banco Santander, Sabadell o Bankinter—, y ajustar sus tablas salariales para permanecer en una posición competitiva. CaixaBank tiene como objetivo declarado situarse en el cuartil superior del mercado bancario español para los perfiles más críticos, con el fin de atraer y retener el talento especializado que el proceso de transformación digital del grupo exige.

\subsection{Equidad individual}

La equidad individual garantiza que, dentro de una misma banda salarial, se reconozcan las diferencias de rendimiento y aportación de cada empleado. CaixaBank materializa este principio a través de un sistema de \textbf{retribución variable ligada al desempeño}, cuya base es la evaluación de rendimiento anual gestionada mediante SAP SuccessFactors.

En función de los resultados obtenidos, el rendimiento de cada empleado se clasifica en distintos niveles de desempeño que determinan el porcentaje de retribución variable al que tiene derecho. De esta forma, un gestor comercial que supere de manera significativa sus objetivos anuales percibirá un bonus considerablemente superior al de un compañero del mismo nivel que haya alcanzado únicamente los mínimos. Este mecanismo busca no solo reconocer el mérito individual, sino también mantener el incentivo y la motivación de los equipos de alto rendimiento.

\section{Retribución por puesto vs.\ por competencias}

Uno de los debates centrales en el diseño de sistemas retributivos es si la base de la compensación debe ser el puesto de trabajo en sí mismo —con independencia de quién lo ocupe— o bien las competencias y capacidades que el empleado aporta a la organización. CaixaBank aplica un \textbf{modelo híbrido} que combina ambos enfoques de forma complementaria.

Por un lado, la \textbf{retribución basada en el puesto} constituye el componente fijo y estable del salario. El convenio colectivo sectorial y la estructura de niveles profesionales internos de la entidad asignan a cada categoría un rango salarial de referencia, determinado por las responsabilidades formales del puesto, su posición en el organigrama y el impacto esperado en los resultados del negocio. Esta base fija garantiza la equidad interna y la seguridad económica del empleado con independencia de las fluctuaciones del rendimiento en un periodo concreto.

Por otro lado, la \textbf{retribución basada en competencias} actúa como elemento diferenciador sobre esa base fija. En CaixaBank, determinadas competencias tienen un reconocimiento retributivo explícito:

\begin{itemize}
    \item La obtención y mantenimiento de \textbf{certificaciones financieras reguladas} —como la certificación MiFID II para gestores de banca personal o la titulación EFPA/EFA para gestores de banca privada— está directamente vinculada a la posibilidad de acceder a puestos de mayor categoría y, por ende, a bandas salariales superiores.
    \item Las \textbf{habilidades digitales avanzadas} —manejo de herramientas de análisis de datos, uso de inteligencia artificial generativa o gestión de proyectos tecnológicos— son cada vez más valoradas en el proceso de promoción interna, especialmente en el contexto de la transformación digital recogida en el Plan Estratégico 2025--2027.
    \item La \textbf{capacidad de liderazgo} y los resultados obtenidos en programas de desarrollo directivo internos influyen en la posibilidad de acceder a complementos salariales reservados a los perfiles de alta potencial (\textit{high potential}).
\end{itemize}

Este modelo mixto responde a la necesidad de combinar la estabilidad y la equidad que ofrece la retribución por puesto con la flexibilidad y el reconocimiento del mérito individual que aporta el enfoque por competencias, lo que resulta especialmente adecuado para una organización tan heterogénea como CaixaBank.

\section{Tipos de incentivos y elementos no financieros}

La política de compensación total de CaixaBank va mucho más allá del salario base. La entidad diseña un paquete retributivo integral que combina incentivos financieros directos e indirectos con un conjunto de beneficios no financieros cuyo objetivo es mejorar la calidad de vida de sus empleados y fortalecer su compromiso con la organización.

\subsection{Incentivos financieros directos: retribución variable}

\begin{itemize}
    \item \textbf{Bonus comercial anual:} Es el principal incentivo financiero para los perfiles de red de oficinas. Se calcula como un porcentaje del salario fijo y se determina en función del grado de cumplimiento de los objetivos individuales y colectivos (captación de clientes, volumen de contratación de productos, calidad del servicio). En años de buen rendimiento, este componente puede suponer entre el 10\,\% y el 30\,\% de la retribución total del gestor.
    \item \textbf{Planes de co-inversión y acciones para directivos:} Para los perfiles de alta dirección, CaixaBank dispone de planes de incentivos a largo plazo basados en la entrega diferida de acciones del banco. Estos programas vinculan parte de la retribución al desempeño del grupo a lo largo de tres o cinco años, alineando los intereses de los directivos con los de los accionistas y favoreciendo la retención del talento directivo.
    \item \textbf{Incentivos por captación y fidelización:} Algunos programas retributivos contemplan primas puntuales por la incorporación de clientes de alto valor o por la renovación de contratos de productos de largo plazo.
\end{itemize}

\subsection{Retribución indirecta y prestaciones}

\begin{itemize}
    \item \textbf{Seguro médico colectivo:} CaixaBank ofrece a todos sus empleados la posibilidad de acceder a un seguro de salud de empresa con condiciones ventajosas, extensible a los familiares directos. Esta prestación tiene un valor percibido muy alto y contribuye significativamente a la satisfacción del empleado.
    \item \textbf{Plan de pensiones de empresa:} La entidad realiza aportaciones periódicas al plan de pensiones de sus empleados como complemento a la jubilación pública, lo que constituye uno de los elementos retributivos indirectos más apreciados, especialmente entre los empleados de mayor antigüedad.
    \item \textbf{Préstamos hipotecarios y personales a tipo preferencial:} Los trabajadores de CaixaBank tienen acceso a productos de financiación con condiciones más favorables que las del mercado, lo que supone un beneficio económico tangible y directo.
    \item \textbf{Ayuda escolar y guardería:} La entidad contempla subsidios o ayudas económicas para el cuidado de hijos menores, en línea con su compromiso con la conciliación de la vida laboral y familiar.
\end{itemize}

\subsection{Elementos no financieros}

Los elementos no financieros de la retribución constituyen un componente esencial del paquete de compensación total de CaixaBank, especialmente relevantes para atraer a las generaciones más jóvenes:

\begin{itemize}
    \item \textbf{Formación y desarrollo profesional:} El acceso a la plataforma Virtaula, a los programas de las escuelas especializadas del banco y la financiación de certificaciones externas (EFPA, CFA, etc.) representa una inversión en el capital humano del empleado que este percibe como parte de su retribución total.
    \item \textbf{Flexibilidad horaria:} CaixaBank dispone de medidas de flexibilidad en la entrada y salida, así como de jornada intensiva en el periodo estival, lo que mejora la conciliación laboral.
    \item \textbf{Reconocimiento y clima laboral:} La entidad dispone de programas de reconocimiento interno —como premios a la innovación o a la excelencia comercial— que refuerzan el sentido de pertenencia y la motivación sin implicar necesariamente un coste económico directo.
    \item \textbf{Voluntariado corporativo:} La vinculación con la Fundación ``la Caixa'' permite a los empleados participar en actividades de responsabilidad social, lo que refuerza el orgullo de pertenencia y la identificación con los valores de la entidad.
\end{itemize}

\section{Tendencias: retribución flexible y salario emocional}

La gestión de la retribución está evolucionando hacia modelos más personalizados y centrados en el bienestar integral del empleado. CaixaBank ha incorporado dos de las tendencias más relevantes en este ámbito: la retribución flexible y el salario emocional.

\subsection{Retribución flexible}

La retribución flexible —también denominada \textit{flex benefits} o salario en especie— consiste en permitir que el empleado decida voluntariamente transformar una parte de su retribución dineraria en una serie de productos o servicios exentos total o parcialmente de tributación en el Impuesto sobre la Renta de las Personas Físicas (IRPF), optimizando así su poder adquisitivo neto sin incrementar el coste laboral para la empresa.

CaixaBank ofrece a sus empleados la posibilidad de acogerse a un plan de retribución flexible que incluye, entre otros, los siguientes productos:

\begin{itemize}
    \item \textbf{Seguro médico:} Los empleados pueden destinar parte de su salario bruto al pago de la prima del seguro médico privado, beneficiándose de la exención fiscal establecida en el artículo 42 de la Ley del IRPF.
    \item \textbf{Transporte:} Bonificaciones para el uso de transporte público colectivo, con el correspondiente ahorro fiscal.
    \item \textbf{Ticket restaurante:} Vales de comida para los días de trabajo presencial, exentos de IRPF hasta los límites legales.
    \item \textbf{Guardería y cheque escolar:} Contribución a los gastos de cuidado de hijos menores de tres años, con beneficio fiscal para el trabajador.
    \item \textbf{Formación:} Determinados programas de formación financiados por la empresa están excluidos de la base imponible del impuesto.
\end{itemize}

Este sistema aporta a los empleados un aumento efectivo de su renta disponible sin modificar el coste total de personal para la entidad, lo que lo convierte en una herramienta muy eficiente para incrementar la satisfacción retributiva. Además, permite a cada trabajador personalizar su paquete de beneficios en función de su situación personal y familiar, respondiendo a la creciente demanda de individualización de las condiciones laborales.

\subsection{Salario emocional}

El concepto de salario emocional hace referencia al conjunto de beneficios no económicos que una organización ofrece a sus empleados y que tienen un impacto directo en su satisfacción, motivación y bienestar, sin que su valor se traduzca directamente en un ingreso monetario. En un contexto de mercado de talento competitivo, el salario emocional se ha convertido en un factor diferenciador decisivo para la atracción y retención de profesionales, especialmente entre las generaciones más jóvenes.

CaixaBank ha desarrollado diversas iniciativas en esta línea, agrupadas bajo su programa de bienestar y conciliación:

\begin{itemize}
    \item \textbf{Teletrabajo estructurado:} Tras la aceleración impuesta por la pandemia, la entidad ha institucionalizado el trabajo en remoto para los puestos elegibles, permitiendo hasta dos días de teletrabajo semanales. Esta medida impacta directamente en la reducción de los tiempos de desplazamiento y en la autonomía del empleado para organizar su jornada.
    \item \textbf{Medidas de conciliación:} CaixaBank dispone de una batería de medidas que incluyen días adicionales de asuntos propios, permisos de maternidad y paternidad ampliados voluntariamente más allá de los mínimos legales, y adaptaciones de jornada para el cuidado de familiares dependientes.
    \item \textbf{Bienestar mental y físico:} La entidad ha ampliado su oferta de bienestar con programas de apoyo psicológico, acceso a plataformas de mindfulness y actividad física, y talleres de gestión del estrés, en respuesta al incremento de las bajas por salud mental en el sector bancario.
    \item \textbf{Voluntariado y propósito:} La participación en los programas sociales de la Fundación ``la Caixa'' —que opera en ámbitos como la inserción laboral, la investigación médica o la educación— proporciona a los empleados un sentido de propósito que trasciende su función laboral diaria, aumentando su vinculación emocional con la organización.
    \item \textbf{Desarrollo y formación como beneficio percibido:} El acceso privilegiado a un ecosistema de formación continua —desde las escuelas internas hasta las certificaciones financieras costeadas por el banco— es percibido por los empleados como una inversión en su futuro profesional que difícilmente encontrarían en otra entidad, lo que actúa como un poderoso elemento de retención.
\end{itemize}

En definitiva, la combinación de retribución flexible y salario emocional responde a la comprensión de que la motivación del empleado moderno no depende únicamente de su nivel salarial, sino de la calidad de su experiencia global como trabajador. CaixaBank, consciente de ello, integra ambas tendencias en su propuesta de valor al empleado (\textit{Employee Value Proposition}) como palanca estratégica para mantener una plantilla comprometida, productiva y alineada con los objetivos del grupo.



% -------------------------------------------------------------------
%  BACKMATTER
%  (bibliografía y fuentes consultadas)
% -------------------------------------------------------------------
\backmatter
\pdfbookmark[-1]{Referencias}{BM-Referencias}

% Fuentes consultadas (URLs y documentos teóricos)
% !TeX root = ../practica.tex
% !TeX encoding = utf8
%
% ====================================================================
%  BIBLIOGRAFÍA Y FUENTES CONSULTADAS T6, 7 y 8
% ====================================================================

\chapter*{Fuentes consultadas (Temas 6, 7 y 8)}
\addcontentsline{toc}{chapter}{Fuentes consultadas (Temas 6, 7 y 8)}

% ================================================================
\section*{Fuentes sobre CaixaBankCareers y procesos de Selección}
% ================================================================

\begin{enumerate}[label={[\arabic*]}]

  \item \textbf{CaixaBankCareers} --- ``Página oficial de la bolsa de empleo e información de etapas de reclutamiento''.\\
        \url{https://caixabankcareers.com/}

  \item \textbf{Jobseeker} --- Análisis de los requisitos para la criba curricular.\\
        \url{https://www.jobseeker.com/es/}

  \item \textbf{Ejemplos-Curriculum} --- Revisión sobre los procesos bancarios automatizados y filtros (\textit{Killer questions}).\\
        \url{https://www.ejemplos-curriculum.com/blog/entrevista-trabajo-la-caixa/}

  \item \textbf{Foros de empleo y Oposiciones} (Infoempleo, Randstad, Buscaoposiciones) --- Experiencias reales de candidatos en dinámicas de grupo y pruebas psicotécnicas de CaixaBank.\\
        \url{https://www.infoempleo.com/}
% Fuentes de jose angel:
  \item \textbf{LinkedIn --- CaixaBank} --- Perfil corporativo; análisis de la estrategia de \textit{employer branding} y publicación de ofertas (reclutamiento 3.0).\\
        \url{https://es.linkedin.com/company/caixabank/jobs}

  \item \textbf{People Xperience Hub / Dialnet} --- Artículo académico sobre la iniciativa RPO de CaixaBank (2019): estandarización y digitalización del ciclo de atracción de talento.\\
        \url{https://dialnet.unirioja.es/servlet/articulo?codigo=7731413}

  \item \textbf{Entrevista con directora de sucursal de CaixaBank} --- Entrevista semiestructurada realizada el 11 de abril de 2026. Fuente primaria sobre el proceso de reclutamiento y selección aplicado en la entidad: métodos utilizados, herramientas digitales, videoentrevistas e iniciativas en curso (IA, gamificación).

\end{enumerate}

% ================================================================
\section*{Documentos teóricos de referencia}
% ================================================================

\begin{enumerate}[label={[\arabic*]}, resume]

  \item Apuntes de clase de los Temas 6, 7 y 8.

  \item \textbf{Guión para el desarrollo de la práctica grupal Temas 6, 7 y 8},
        Dirección de Recursos Humanos~I, Doble Grado ADE-Ingenierías,
        Universidad de Granada, curso 2025--2026.

\end{enumerate}

\endinput


% Bibliografía BibTeX (si se añaden citas con \cite{})
% \bibliographystyle{alpha-es}
% \begin{small}
%   \bibliography{library.bib}
% \end{small}

\end{document}
